\documentclass[../main.tex]{subfiles}
\begin{document}
% SOURCE p=226; printed=211; object=ch07-title; region=top
\chapter{Radon 测度}\label{chapter:ch07-radon-measures}

% SOURCE p=226; printed=211; object=ch07-introduction; region=middle
本章研究局部紧 Hausdorff（LCH）空间上的测度与积分理论。我们在\cref{section:ch02-n-dimensional-lebesgue-integral}中已经看到，$\mathbb R^n$ 上的 Lebesgue 测度与其拓扑配合良好：可测集可以用开集或紧集逼近，可积函数可以用连续函数逼近。在更一般的空间上研究具有类似性质的测度，是很自然的问题。此外，连续函数空间上的某些线性泛函，恰可表示为关于这类测度的积分。这个事实在测度论与泛函分析之间建立了重要联系，也提供了一种构造测度的有力工具。

% SOURCE p=226; printed=211; object=ch07-standing-conventions; region=middle-lower
\begin{definition}[本章的空间与 Borel 记号]\label{definition:ch07-borel-conventions}
本章始终以 $X$ 表示 LCH 空间。继续使用\cref{chapter:ch01-measures}中在度量空间背景下引入的术语：$\mathscr B_X$ 表示由 $X$ 的开集生成的 \textbf{Borel $\sigma$ 代数}，其上的测度称为 \textbf{Borel 测度}；闭集的可数并称为 $F_\sigma$ 集，开集的可数交称为 $G_\delta$ 集。
\end{definition}
% SUPPLEMENT: id=v14-ch07-borel-conventions; parent=ch07-standing-conventions; kind=motivation+hypotheses; origin=independent-derivation
\begin{remark}
\textbf{动机与逐条件说明（重构新增）。}Borel 结构把拓扑给出的开集纳入测度定义域，使“从外以开集逼近”和“从内以紧集逼近”有意义。Hausdorff 性保证紧集闭，因而紧集是 Borel 集；局部紧性则提供紧邻域及紧支撑的连续截断函数。这里无需假设空间可度量或有可数基；缺乏可数性时，任意开覆盖未必可化成可数覆盖，所以不能把 $\mathbb R^n$ 上所有正则性论证直接照搬。除另作说明外，本章的 Borel 测度取非负值；复测度将在相应部分另行指定。$F_\sigma$ 与 $G_\delta$ 中的“可数”是固定定义，不能替换为任意并或交。
\end{remark}
% END SUPPLEMENT: id=v14-ch07-borel-conventions
\subfile{Chapter-07/section-01.tex}
\subfile{Chapter-07/section-02.tex}
\subfile{Chapter-07/section-03.tex}
\subfile{Chapter-07/section-04.tex}
\subfile{Chapter-07/section-05.tex}
\end{document}
